% !TeX root = main.tex

% Изображения добавляются в images; до этого вместо них выводятся рамки.
\newcommand{\reportimage}[3]{%
  \begin{figure}[htbp]
    \centering
    \IfFileExists{images/#1}{%
      \includegraphics[width=0.95\textwidth]{#1}%
    }{%
      \fbox{%
        \parbox[c][3cm][c]{0.90\textwidth}{%
          \centering
          Место для рисунка\\[0.4em]
          \texttt{images/\detokenize{#1}}
        }%
      }%
    }
    \caption{#2}
    \label{#3}
  \end{figure}
}

\tableofcontents
\newpage

\section{Цель работы}

Освоить процедуру синтеза адаптивного наблюдателя состояния линейного
объекта. Исследовать влияние спектрального состава входного сигнала
и коэффициента адаптации на восстановление состояния и оценивание
неизвестных параметров.

\section{Постановка задачи и исходные данные}

Рассматривается асимптотически устойчивый объект второго порядка
\begin{equation}
  \dot{\bm x}=A\bm x+\bm b u,
  \qquad y=C\bm x,
  \qquad \bm x\in\mathbb R^2,
  \label{eq:plant}
\end{equation}
где
\begin{equation}
  A=\begin{pmatrix}-a_1&1\\-a_0&0\end{pmatrix},
  \qquad
  \bm b=\begin{pmatrix}b_1\\b_0\end{pmatrix},
  \qquad C=\begin{pmatrix}1&0\end{pmatrix}.
  \label{eq:plant-structure}
\end{equation}
Измерению доступны вход $u(t)$ и выход $y(t)=x_1(t)$.
Коэффициенты $a_0,a_1,b_0,b_1$ считаются неизвестными наблюдателю.
Необходимо построить оценки параметров и вектора состояния
$\hat{\bm x}(t)$, обеспечивающие при выполнении условий сходимости
\begin{equation}
  \lim_{t\to\infty}\norm{\bm x(t)-\hat{\bm x}(t)}_2=0.
  \label{eq:goal}
\end{equation}


Исходные данные:
\begin{equation}
  a_1=1,\qquad a_0=1,\qquad b_1=1,\qquad b_0=2,
  \qquad k_1=0{,}2,\qquad k_0=0{,}01.
  \label{eq:variant}
\end{equation}
Следовательно,
\begin{equation}
  A=\begin{pmatrix}-1&1\\-1&0\end{pmatrix},
  \qquad
  \bm b=\begin{pmatrix}1\\2\end{pmatrix},
  \qquad C=\begin{pmatrix}1&0\end{pmatrix}.
\end{equation}
Передаточная функция объекта при нулевых начальных условиях имеет вид
\begin{equation}
  W(s)=C(sI_2-A)^{-1}\bm b
  =\frac{b_1s+b_0}{s^2+a_1s+a_0}
  =\frac{s+2}{s^2+s+1}.
\end{equation}
Собственные значения матрицы $A$ равны
$-1/2\pm \mathrm{i}\sqrt{3}/2$, поэтому объект асимптотически устойчив.

Для построения наблюдателя используется известная вспомогательная матрица
\begin{equation}
  A_0=\begin{pmatrix}-k_1&1\\-k_0&0\end{pmatrix}
  =\begin{pmatrix}-0{,}2&1\\-0{,}01&0\end{pmatrix}.
  \label{eq:filter-matrix}
\end{equation}
Её характеристический полином
\begin{equation}
  K(s)=s^2+k_1s+k_0=s^2+0{,}2s+0{,}01=(s+0{,}1)^2
  \label{eq:filter-polynomial}
\end{equation}
гурвицев: оба корня равны $-0{,}1$.

Истинные коэффициенты объекта используются для его моделирования
и вычисления ошибок оценивания. В алгоритм наблюдателя передаются
только измеряемые сигналы $u,y$, известные коэффициенты $k_0,k_1$,
коэффициент адаптации и начальные оценки.

\section{Синтез адаптивного наблюдателя}

\subsection{Параметризованное представление объекта}

Благодаря канонической форме матриц выполняется
\begin{equation}
  A-A_0=\bm d C,
  \qquad
  \bm d=\begin{pmatrix}k_1-a_1\\k_0-a_0\end{pmatrix},
  \qquad
  \dot{\bm x}=A_0\bm x+\bm d y+\bm b u.
  \label{eq:matching}
\end{equation}
Вектор оцениваемых параметров в соответствии с (5.4) и (5.7)
методических указаний:
\begin{equation}
  \bm\theta=
  \begin{pmatrix}
    k_0-a_0\\k_1-a_1\\b_0\\b_1
  \end{pmatrix}
  =
  \begin{pmatrix}
    -0{,}99\\-0{,}8\\2\\1
  \end{pmatrix}.
  \label{eq:true-parameters}
\end{equation}

Введём фильтрованные сигналы
\begin{equation}
  \xi_0=\frac{1}{K(s)}[y],\qquad
  \xi_1=\frac{s}{K(s)}[y],\qquad
  v_0=\frac{1}{K(s)}[u],\qquad
  v_1=\frac{s}{K(s)}[u].
  \label{eq:filtered-signals}
\end{equation}
Реализация фильтров:
\begin{equation}
  \begin{aligned}
    \dot\xi_0&=\xi_1,
    &\dot\xi_1&=-k_0\xi_0-k_1\xi_1+y,\\
    \dot v_0&=v_1,
    &\dot v_1&=-k_0v_0-k_1v_1+u.
  \end{aligned}
  \label{eq:filter-equations}
\end{equation}
Начальные состояния фильтров принимаются нулевыми.

При нулевом начальном состоянии объекта его выход представляется
линейной регрессией
\begin{equation}
  y=\bm\theta^{\mathsf T}\bm\omega,
  \qquad
  \bm\omega=
  \begin{pmatrix}\xi_0&\xi_1&v_0&v_1\end{pmatrix}^{\mathsf T}.
  \label{eq:output-regression}
\end{equation}
\subsection{Настраиваемая модель и алгоритм адаптации}

Настраиваемая модель и ошибки оценивания:
\begin{equation}
  \hat y=\hat{\bm\theta}^{\mathsf T}\bm\omega,
  \qquad
  \varepsilon=y-\hat y,
  \qquad
  \tilde{\bm\theta}=\bm\theta-\hat{\bm\theta}.
  \label{eq:adjustable-model}
\end{equation}
Градиентный алгоритм адаптации:
\begin{equation}
  \dot{\hat{\bm\theta}}=\gamma\bm\omega\varepsilon,
  \qquad \gamma>0.
  \label{eq:adaptation}
\end{equation}
Коэффициент $\gamma$ выбирается экспериментально.
Оценки коэффициентов объекта:
\begin{equation}
  \hat a_0=k_0-\hat\theta_1,\qquad
  \hat a_1=k_1-\hat\theta_2,\qquad
  \hat b_0=\hat\theta_3,\qquad
  \hat b_1=\hat\theta_4.
\end{equation}

В модели без начального переходного слагаемого
$\varepsilon=\tilde{\bm\theta}^{\mathsf T}\bm\omega$.
Для функции Ляпунова
\begin{equation}
  V=\frac{1}{2\gamma}
  \tilde{\bm\theta}^{\mathsf T}\tilde{\bm\theta}
\end{equation}
с учётом постоянства истинных параметров получаем
\begin{equation}
  \dot V
  =-\frac{1}{\gamma}
    \tilde{\bm\theta}^{\mathsf T}\dot{\hat{\bm\theta}}
  =-\tilde{\bm\theta}^{\mathsf T}\bm\omega\varepsilon
  =-\varepsilon^2\le0.
  \label{eq:lyapunov-derivative}
\end{equation}
При ограниченности $\bm\omega$ и $\dot{\bm\omega}$ отсюда следует
асимптотическое затухание $\varepsilon$ в рассматриваемой модели.
Неположительность $\dot V$ сама по себе не гарантирует
сходимости всех оценок параметров к истинным значениям.

Достаточным условием сходимости параметрических ошибок является
неисчезающее возбуждение регрессора: существуют $T>0$ и $\alpha>0$,
для которых
\begin{equation}
  \int_t^{t+T}\bm\omega(\tau)\bm\omega^{\mathsf T}(\tau)\,d\tau
  \succeq\alpha I_4
  \qquad\text{для всех }t\ge t_0.
  \label{eq:pe}
\end{equation}
\begin{samepage}
\subsection{Восстановление вектора состояния}

Для $n=2$, $m=1$ и порядка параметров~\eqref{eq:true-parameters}
формула (6.6) методических указаний принимает вид
\begin{equation}
  \begin{aligned}
    \hat{\bm x}={}&
    \hat\theta_1\Phi(s)\bm e_2[y]
    +\hat\theta_2\Phi(s)\bm e_1[y]\\
    &+\hat\theta_3\Phi(s)\bm e_2[u]
    +\hat\theta_4\Phi(s)\bm e_1[u],
  \end{aligned}
  \label{eq:state-estimate-expanded}
\end{equation}
\end{samepage}
где $\bm e_1=(1,0)^{\mathsf T}$, $\bm e_2=(0,1)^{\mathsf T}$.
Для матрицы~\eqref{eq:filter-matrix} получаем
\begin{equation}
  \Phi(s)=(sI_2-A_0)^{-1}
  =
  \begin{pmatrix}s+k_1&-1\\k_0&s\end{pmatrix}^{-1}
  =
  \frac{1}{K(s)}
  \begin{pmatrix}s&1\\-k_0&s+k_1\end{pmatrix}.
  \label{eq:filter-resolvent}
\end{equation}
С учётом определений~\eqref{eq:filtered-signals} векторные сигналы
в~\eqref{eq:state-estimate-expanded} имеют вид
\begin{equation}
  \begin{aligned}
    \Phi(s)\bm e_2[y]
      &=\begin{pmatrix}\xi_0\\\xi_1+k_1\xi_0\end{pmatrix},
    &\Phi(s)\bm e_1[y]
      &=\begin{pmatrix}\xi_1\\-k_0\xi_0\end{pmatrix},\\[0.6em]
    \Phi(s)\bm e_2[u]
      &=\begin{pmatrix}v_0\\v_1+k_1v_0\end{pmatrix},
    &\Phi(s)\bm e_1[u]
      &=\begin{pmatrix}v_1\\-k_0v_0\end{pmatrix}.
  \end{aligned}
  \label{eq:vector-filter-signals}
\end{equation}
Объединяя векторные сигналы~\eqref{eq:vector-filter-signals}
в столбцы матрицы $R(t)$ в порядке множителей
$\hat\theta_1,\ldots,\hat\theta_4$, получаем
\begin{equation}
  \begin{gathered}
    \hat{\bm x}=R(t)\hat{\bm\theta},
    \qquad \hat x_1=\hat y,\\[0.4em]
    R(t)=
    \begin{pmatrix}
      \xi_0&\xi_1&v_0&v_1\\
      \xi_1+k_1\xi_0&-k_0\xi_0&
      v_1+k_1v_0&-k_0v_0
    \end{pmatrix}.
  \end{gathered}
  \label{eq:state-estimate}
\end{equation}


\section{Порядок моделирования}

Структура модели в Simulink (рисунок~\ref{fig:observer-scheme})
включает объект~\eqref{eq:plant}, фильтры~\eqref{eq:filter-equations},
настраиваемую модель~\eqref{eq:adjustable-model}, алгоритм адаптации
\eqref{eq:adaptation} и оценку состояния~\eqref{eq:state-estimate}.

\reportimage{adaptive_observer_scheme.png}
  {Схема моделирования объекта и адаптивного наблюдателя}
  {fig:observer-scheme}

Настройки расчёта приведены в таблице~\ref{tab:simulation-settings}.
% Заполнить таблицу после выбора начальных условий и настроек решателя.
\begin{table}[htbp]
  \centering
  \caption{Настройки моделирования}
  \label{tab:simulation-settings}
  \begin{tabularx}{\textwidth}{@{}X l@{}}
    \toprule
    Параметр & Значение\\
    \midrule
    Начальное состояние объекта $\bm x(0)$
      & \underline{\hspace{3cm}}\\
    Начальная оценка параметров $\hat{\bm\theta}(0)$
      & \underline{\hspace{3cm}}\\
    Начальные состояния фильтров & Нулевые\\
    Время моделирования $t_{\mathrm{mod}}$
      & \underline{\hspace{3cm}}\\
    Решатель & \underline{\hspace{3cm}}\\
    Шаг или ограничения шага интегрирования
      & \underline{\hspace{3cm}}\\
    Допуски решателя & \underline{\hspace{3cm}}\\
    \bottomrule
  \end{tabularx}
\end{table}

В соответствии с заданием выполняются два эксперимента.
В каждом регистрируются норма ошибки состояния
$\norm{\bm x-\hat{\bm x}}_2$ и четыре компоненты
параметрической ошибки $\tilde{\bm\theta}$.
Для сравнения влияния входных сигналов начальные условия
и настройки решателя сохраняются одинаковыми.

\clearpage
\section{Эксперимент с одногармоническим входным сигналом}

В первом эксперименте задаётся вход
\begin{equation}
  u_1(t)=10\sin t.
\end{equation}
Коэффициент адаптации выбирается экспериментально:
$\gamma_1=\underline{\hspace{2cm}}$.
Норма ошибки состояния: рисунок~\ref{fig:single-state-error};
параметрические ошибки: рисунок~\ref{fig:single-parameter-errors}.

\reportimage{single_state_error.png}
  {Норма ошибки состояния $\norm{\bm x-\hat{\bm x}}_2$ при $u_1(t)=10\sin t$}
  {fig:single-state-error}
\reportimage{single_parameter_errors.png}
  {Параметрические ошибки $\tilde\theta_1,\ldots,\tilde\theta_4$
   при одногармоническом входном сигнале}
  {fig:single-parameter-errors}

\subsection{Анализ результатов}

\textit{Раздел заполняется после проведения моделирования.}
% Описать поведение нормы ошибки состояния и каждой параметрической ошибки.
% Указать выбранное gamma и длительность переходного процесса.
% Отдельно оценить восстановление состояния и идентификацию параметров.

\clearpage
\section{Эксперимент с многогармоническим входным сигналом}

Во втором эксперименте используется вход
\begin{equation}
  u_2(t)=10\sin t+5\cos 2t+4\cos 4t+3\cos 8t.
\end{equation}
Коэффициент адаптации:
$\gamma_2=\underline{\hspace{2cm}}$.
Норма ошибки состояния: рисунок~\ref{fig:multi-state-error};
параметрические ошибки: рисунок~\ref{fig:multi-parameter-errors}.

\reportimage{multi_state_error.png}
  {Норма ошибки состояния $\norm{\bm x-\hat{\bm x}}_2$
   при многогармоническом входном сигнале}
  {fig:multi-state-error}
\reportimage{multi_parameter_errors.png}
  {Параметрические ошибки $\tilde\theta_1,\ldots,\tilde\theta_4$
   при многогармоническом входном сигнале}
  {fig:multi-parameter-errors}

\subsection{Анализ результатов}

\textit{Раздел заполняется после проведения моделирования.}
% Сопоставить ошибки состояния и параметров с первым экспериментом.
% Если gamma различаются, учитывать это при сравнении переходных процессов.
% Связать наблюдаемые результаты с возбуждением регрессора,
% не подменяя проверку сходимости только малостью ошибки выхода.

\section{Выводы}

\textit{Выводы формулируются после анализа результатов двух экспериментов.}
% Указать, выполнена ли цель наблюдения и как ведут себя оценки параметров.
% Описать влияние состава входного сигнала и выбранного gamma.
